\textbf{1. Kiến thức cần nhớ}
\begin{itemize}
\item \textbf{Bảng thống kê}: Bảng số liệu phản ánh sự tương ứng giữa đối tượng và số lượng/tỉ lệ.
\item \textbf{Các loại biểu đồ thường dùng}:
\begin{itemize}
\item \textbf{Biểu đồ cột}: Thích hợp so sánh số liệu giữa các đối tượng độc lập.
\item \textbf{Biểu đồ cột kép}: Thích hợp so sánh từng cặp số liệu cùng loại của hai nhóm đối tượng.
\item \textbf{Biểu đồ đoạn thẳng}: Thích hợp biểu diễn \textbf{xu hướng thay đổi} của một số liệu theo thời gian (nhiệt độ theo giờ, dân số qua các năm\dots).
\item \textbf{Biểu đồ hình quạt tròn}: Thích hợp biểu diễn \textbf{tỉ lệ phần trăm} của từng phần so với toàn thể (tổng các phần tử luôn bằng \%$).
\end{itemize}
\end{itemize}

\textbf{2. Phương pháp giải các dạng}

%% DANG 8D5?19-1
\textbf{Dạng 1. Đọc, lập bảng thống kê}
\begin{enumerate}
\item Xác định các hàng và cột thể hiện các tiêu chí cần thống kê.
\item Điền đúng số liệu thu thập được vào các ô tương ứng.
\item Tính tổng các dòng, cột khi cần thiết để kiểm tra tính đầy đủ của số liệu.
\end{enumerate}

%% DANG 8D5?19-2
\textbf{Dạng 2. Chọn và vẽ biểu đồ phù hợp}
\begin{enumerate}
\item Căn cứ vào mục đích biểu diễn để chọn loại biểu đồ:
\begin{itemize}
\item Muốn thấy sự tăng/giảm theo thời gian $\rightarrow$ chọn biểu đồ đoạn thẳng.
\item Muốn thấy cơ cấu phần trăm so với tổng thể $\rightarrow$ chọn biểu đồ hình quạt tròn.
\item Muốn so sánh độ lớn giữa các đối tượng $\rightarrow$ chọn biểu đồ cột (hoặc cột kép).
\end{itemize}
\item Vẽ đầy đủ các thành phần của biểu đồ: Tên biểu đồ, các trục tọa độ, chia khoảng cách tỉ lệ hợp lí, ghi chú đơn vị và chú giải.
\end{enumerate}

%% DANG 8D5?19-3
\textbf{Dạng 3. Chuyển dữ liệu giữa các dạng biểu diễn}
\begin{enumerate}
\item Đọc chính xác giá trị từ biểu đồ (cột, đoạn thẳng hoặc quạt tròn).
\item Lập bảng số liệu tương ứng.
\item Ngược lại, từ bảng số liệu tính tỉ lệ phần trăm để chuyển sang biểu đồ quạt tròn:
80830\text{Tỉ lệ \%} = \dfrac{\text{Giá trị thành phần}}{\text{Tổng thể}} \times 100\%80830
\end{enumerate}
